\subsection{The Corporation (2003)}

El documental analiza el papel de la corporación como una de las instituciones económicas más influyentes del mundo contemporáneo. A través de entrevistas, ejemplos históricos y casos empresariales, la película plantea que las corporaciones están estructuralmente diseñadas para perseguir un objetivo central: \textbf{la maximización del capital y el beneficio para sus accionistas.} Desde esta perspectiva, la película no critica necesariamente la existencia de las corporaciones en sí mismas. Como instituciones económicas, constituyen mecanismos eficientes para organizar inversión, coordinar grandes proyectos productivos y movilizar capital. El problema que el documental pone en evidencia es que, al estar guiadas principalmente por incentivos financieros, las decisiones corporativas suelen evaluarse únicamente mediante métricas monetarias, dejando en segundo plano otras dimensiones sociales o humanas.

Este enfoque permite comprender por qué muchas de las prácticas que aparecen en el documental no surgen necesariamente de una intención maliciosa individual, sino de la lógica interna del sistema. Si una corporación está diseñada para priorizar la rentabilidad y el crecimiento del capital, entonces sus decisiones tenderán a optimizar esos resultados, incluso cuando las consecuencias externas puedan ser negativas. En este sentido, el comportamiento de los actores dentro del sistema responde a incentivos racionales desde el punto de vista económico.

Un ejemplo particularmente inquietante es el caso de un trader que expresa entusiasmo al observar que los ataques del 11 de septiembre de 2001 provocaron un aumento en el precio del oro, activo en el que él se encontraba invertido. Esta escena resulta perturbadora no tanto por la emoción inicial en sí misma, sino por lo que revela sobre la relación entre los mercados financieros y los eventos humanos. Antes de ser un gestor de portafolio, el trader es una persona que reacciona ante los resultados de su inversión; desde esa perspectiva, no es necesariamente incorrecto experimentar satisfacción cuando una posición financiera genera ganancias inesperadas. Sin embargo, la cuestión ética emerge al considerar cómo se procesa y se responde a esa situación. Cuando una tragedia humana se convierte únicamente en una oportunidad financiera, el riesgo es que los incentivos del sistema fomenten una separación cada vez mayor entre las consecuencias humanas de los eventos y las reacciones de los agentes económicos. La crítica que puede derivarse de esta escena no es que los participantes del mercado obtengan beneficios de cambios inesperados en los precios, sino que la lógica institucional puede normalizar respuestas que ignoran la dimensión humana de esos acontecimientos.

El documental sugiere que el problema no reside únicamente en las acciones individuales, sino en la estructura de incentivos que orienta el comportamiento corporativo. Cuando la maximización del capital se convierte en el criterio dominante de decisión, prácticas como la explotación de crisis económicas, la externalización de riesgos o la reducción abrupta de costos laborales pueden aparecer como decisiones racionales dentro del marco corporativo. La crítica central, por tanto, no es la existencia de las corporaciones, sino el hecho de que su diseño institucional privilegia de manera casi exclusiva la acumulación de capital, lo que puede generar tensiones profundas entre los objetivos económicos de las organizaciones y las implicaciones sociales de sus decisiones.
